\chapter{分块、Packing 与微内核}

上一章看到，改变循环顺序可以改善一部分访问模式，但还不够。真正高性能的 GEMM 往往会分成多层结构：先把矩阵切成块，再把块复制成更适合内核访问的 packed 格式，最后用一个手写或生成的微内核在寄存器里计算小块 $C$。这一章解释这些结构为什么存在。

\section{分块的根本目的}

分块的目标是控制工作集大小。假设我们计算 $C$ 的一个小块，它需要 $A$ 的若干行和 $B$ 的若干列。只要这部分 $A$、$B$、$C$ 能留在合适缓存层级里，内层计算就能反复复用它们。若块太大，缓存装不下，数据被提前挤出；若块太小，循环开销、边界处理和 packing 成本占比过高。

可以先写一个教学版分块，不考虑最优块大小。

\begin{lstlisting}[language=C++,caption={教学版分块 GEMM：按块组织复用}]
#include <algorithm>
#include <cstdint>
#include <span>

void gemm_blocked(std::span<const float> a,
                  std::span<const float> b,
                  std::span<float> c,
                  std::int32_t m,
                  std::int32_t n,
                  std::int32_t k_size,
                  std::int32_t block) {
    std::fill(c.begin(), c.end(), 0.0F);
    for (std::int32_t ii = 0; ii < m; ii += block) {
        for (std::int32_t kk = 0; kk < k_size; kk += block) {
            for (std::int32_t jj = 0; jj < n; jj += block) {
                const std::int32_t i_end = std::min(ii + block, m);
                const std::int32_t k_end = std::min(kk + block, k_size);
                const std::int32_t j_end = std::min(jj + block, n);
                for (std::int32_t i = ii; i < i_end; ++i) {
                    for (std::int32_t k = kk; k < k_end; ++k) {
                        const float av =
                            a[static_cast<std::size_t>(i) * k_size + k];
                        for (std::int32_t j = jj; j < j_end; ++j) {
                            c[static_cast<std::size_t>(i) * n + j] +=
                                av * b[static_cast<std::size_t>(k) * n + j];
                        }
                    }
                }
            }
        }
    }
}
\end{lstlisting}

这段代码还不是工业实现，但它已经表达了分块思想。它不是减少总 FLOPs，而是改变数据在缓存里的停留方式。分块参数需要按目标机器调优，因为 L1、L2、L3 容量、cache associativity、SIMD 宽度、寄存器数量和预取行为都会影响最佳值。

\section{为什么需要 Packing}

分块之后，仍然可能存在跨步访问和边界复杂性。Packing 的做法是把即将计算的一块 $A$ 或 $B$ 复制到连续、对齐、顺序访问的临时缓冲区。初看这像额外成本：明明原矩阵里已经有数据，为什么还要复制一次？答案是复制一次可以换来内层微内核大量规则访问。只要 packed block 被复用足够多次，复制成本就能被摊薄。

Packing 还让微内核更简单。微内核不必处理原矩阵的 leading dimension、跨步、边界和奇怪布局，只要按固定格式读取连续数据。许多库会把 $B$ 的 panel 打包成适合连续向量加载的形状，把 $A$ 的小块打包成适合广播或向量加载的形状。这样内层循环就能专注于 FMA 和寄存器调度。

\begin{mentalmodel}
Packing 是一次主动的数据重排。它用可控的顺序拷贝，换取后续大量计算中的规则加载、少分支、少地址计算和更稳定的预取效果。
\end{mentalmodel}

\section{微内核和寄存器块}

微内核通常计算 $C$ 的一个很小 tile，例如 $M_R \times N_R$。这个 tile 的累加值保存在寄存器中，循环遍历 $K$ 维，每一步从 $A$ 和 $B$ 取数据并执行 FMA。等 $K$ 维完成后，再把寄存器里的 tile 写回 $C$。

为什么要把 $C$ 的小块放在寄存器里？因为 $C$ 的每个元素要累加很多次。如果每次乘加都读写内存，流量巨大；若把部分和保存在寄存器，内层只在最后写回一次。寄存器数量决定了 tile 不能无限大，SIMD 宽度决定了 $N_R$ 或 $M_R$ 常常按向量 lane 组织。

\begin{lstlisting}[language=C++,caption={教学化微内核形状：寄存器保存 2x4 的 C 小块}]
#include <cstdint>
#include <span>

void microkernel_2x4(std::span<const float> a,
                     std::span<const float> b,
                     std::span<float> c,
                     std::int32_t k_size,
                     std::int32_t c_stride) {
    float c00 = 0.0F;
    float c01 = 0.0F;
    float c02 = 0.0F;
    float c03 = 0.0F;
    float c10 = 0.0F;
    float c11 = 0.0F;
    float c12 = 0.0F;
    float c13 = 0.0F;

    for (std::int32_t k = 0; k < k_size; ++k) {
        const float a0 = a[static_cast<std::size_t>(k)];
        const float a1 = a[static_cast<std::size_t>(k_size) + k];
        const std::size_t b_base = static_cast<std::size_t>(k) * 4U;
        const float b0 = b[b_base + 0U];
        const float b1 = b[b_base + 1U];
        const float b2 = b[b_base + 2U];
        const float b3 = b[b_base + 3U];
        c00 += a0 * b0;
        c01 += a0 * b1;
        c02 += a0 * b2;
        c03 += a0 * b3;
        c10 += a1 * b0;
        c11 += a1 * b1;
        c12 += a1 * b2;
        c13 += a1 * b3;
    }

    c[0] = c00;
    c[1] = c01;
    c[2] = c02;
    c[3] = c03;
    c[static_cast<std::size_t>(c_stride) + 0U] = c10;
    c[static_cast<std::size_t>(c_stride) + 1U] = c11;
    c[static_cast<std::size_t>(c_stride) + 2U] = c12;
    c[static_cast<std::size_t>(c_stride) + 3U] = c13;
}
\end{lstlisting}

这段代码没有使用 intrinsic，但已经体现微内核思想。真正的 SIMD 版本会把一行或多行 $B$ 装进向量寄存器，把 $A$ 的标量广播到向量，执行向量 FMA。寄存器块越大，复用越充分，但寄存器压力也越大。超过寄存器容量后，编译器会 spill 到栈上，性能反而下降。

\section{工业库的分层视角}

OpenBLAS、BLIS、oneDNN 和其他库的 GEMM 结构看起来复杂，是因为它们要同时服务不同矩阵形状、不同数据类型、不同指令集和不同缓存层级。常见结构可以粗略理解为：外层按大块切分，适配 L3 和多线程；中层 packing panel，适配 L2；内层微内核，适配 L1、寄存器和 SIMD。源码阅读时不要一上来陷入所有宏和平台分支，先找出这三层结构。

\begin{exercisebox}
阅读一个开源 BLAS 的 GEMM 路径，记录它在哪里选择 block size，在哪里 packing，在哪里调用架构相关 microkernel。不要求读懂所有汇编，先把数据从原矩阵到 packed buffer 再到 microkernel 的路径画出来。
\end{exercisebox}
